\section{Ассоциативные контейнеры, компараторы, сравнение и сортировка}

\subsection{Каково назначение параметра \texttt{Compare} в шаблонах \texttt{std::set<Key, Compare>} и \texttt{std::map<Key, Value, Compare>}?}
Параметр \texttt{Compare} задаёт правило упорядочивания элементов внутри ассоциативного контейнера. Именно через него контейнер понимает, какой из двух ключей «меньше», а значит -- куда в дереве нужно спускаться при поиске, вставке и удалении.

\begin{definition}[Компаратор]
Компаратор -- это объект-функция (function object), который по двум значениям возвращает \texttt{true}, если первое должно идти раньше второго в заданном порядке.
\end{definition}

Для \texttt{std::set<Key, Compare>} компаратор сравнивает сами элементы типа \texttt{Key}.  
Для \texttt{std::map<Key, Value, Compare>} компаратор сравнивает только ключи \texttt{Key}, а не пары целиком.

\begin{remark}
В \texttt{std::map} порядок элементов определяется именно по ключу, потому что ключ однозначно задаёт место элемента в дереве.
\end{remark}

\subsection{Какая типичная реализация у \texttt{std::set} и \texttt{std::map}? Почему при поиске используется сравнение?}
Типичная реализация ассоциативных контейнеров -- это \textbf{сбалансированное бинарное дерево поиска} (обычно красно-чёрное дерево). Поэтому для определения пути к элементу нужно сравнивать ключи: идти в левое или правое поддерево.

При использовании компаратора по умолчанию таким сравнением является операция ``меньше'', то есть \texttt{operator<}. Именно поэтому ключи должны быть \textit{comparable}: их можно упорядочить некоторым строгим слабым порядком (\textit{strict weak ordering}).

Отсюда следуют типичные асимптотики:
\[
\text{find, insert, erase} = O(\log n).
\]

\subsection{Какие стандартные компараторы используются по умолчанию?}
По умолчанию в упорядоченных ассоциативных контейнерах используется \texttt{std::less<T>}.

\begin{definition}[\texttt{std::less<T>}]
\texttt{std::less<T>} -- компаратор, который реализует сравнение вида \texttt{a < b}.
\end{definition}

Аналогично, \texttt{std::greater<T>} задаёт обратный порядок, то есть сравнение вида \texttt{a > b}.

\begin{itemize}
    \item \texttt{std::set<int>} по умолчанию отсортирован по возрастанию.
    \item \texttt{std::set<int, std::greater<int>>} -- по убыванию.
\end{itemize}

При использовании \texttt{std::less<T>} тип \texttt{T} должен быть сравним через \texttt{operator<}. Если такого оператора нет, то нужно либо определить его, либо передать пользовательский компаратор.

\subsection{Как задаётся пользовательский компаратор?}
Пользовательский компаратор можно задать несколькими способами:
\begin{itemize}
    \item через структуру или класс с перегруженным \texttt{operator()};
    \item через лямбда-выражение, если оно подходит по типу;
    \item через готовый функциональный объект из стандартной библиотеки.
\end{itemize}

\begin{example}
\begin{verbatim}
struct ByAbs {
    bool operator()(int a, int b) const {
        return std::abs(a) < std::abs(b);
    }
};

std::set<int, ByAbs> s;
\end{verbatim}
\end{example}

Здесь множество упорядочивается не по обычному значению, а по модулю.

\begin{remark}
Компаратор должен задавать корректный строгий слабый порядок. Если это свойство нарушить, контейнер может работать некорректно.
\end{remark}

\subsection{Каковы способы реализации пользовательского компаратора?}
На практике чаще всего используют такие формы:
\begin{enumerate}
    \item \textbf{Структура-функтор}:
\begin{verbatim}
struct Cmp {
    bool operator()(const T& a, const T& b) const {
        ...
    }
};
\end{verbatim}

    \item \textbf{Лямбда}, если компаратор нужен локально:
\begin{verbatim}
auto cmp = [](const T& a, const T& b) {
    return ...;
};
\end{verbatim}

    \item \textbf{std::function}, если нужен полиморфизм, но обычно это избыточно и медленнее.
\end{enumerate}

Для ассоциативных контейнеров особенно удобен функтор, потому что тип компаратора должен быть известен на этапе компиляции.

\subsection{Чем пользовательский компаратор отличается от перегрузки \texttt{operator<}?}
Оба подхода решают задачу упорядочивания, но делают это по-разному.

\begin{itemize}
    \item \textbf{Перегрузка \texttt{operator<}} задаёт естественный порядок самого типа.
    \item \textbf{Пользовательский компаратор} задаёт порядок только для конкретного контейнера или алгоритма.
\end{itemize}

Преимущества пользовательского компаратора:
\begin{itemize}
    \item можно хранить один и тот же тип в разных контейнерах с разными порядками;
    \item не нужно навязывать всем пользователям класса единственный вариант сравнения;
    \item удобно для случаев, когда порядок зависит от задачи, а не от природы типа.
\end{itemize}

\begin{example}
Один и тот же тип \texttt{Person} можно отсортировать:
\begin{itemize}
    \item по фамилии в одном контейнере;
    \item по возрасту в другом;
    \item по id в третьем.
\end{itemize}
Если бы порядок задавался только через \texttt{operator<}, это было бы менее гибко.
\end{example}

\subsection{Как устроена «автоматическая сортировка» ассоциативных контейнеров?}
Упорядоченные ассоциативные контейнеры сами поддерживают порядок при каждой вставке. Это не отдельная операция сортировки, а свойство внутренней структуры данных.

То есть:
\begin{itemize}
    \item при \texttt{insert} элемент помещается в нужную позицию;
    \item при \texttt{erase} дерево перестраивается;
    \item при \texttt{find} путь ищется по сравнению ключей;
    \item при обходе контейнера элементы уже идут в отсортированном порядке.
\end{itemize}

\begin{remark}
У \texttt{std::set} и \texttt{std::map} нельзя сказать: «отсортируй контейнер ещё раз». Они уже всегда хранят элементы в отсортированном виде.
\end{remark}

\subsection{Как «сортируют» элементы непустого \texttt{std::set}?}
Формально \texttt{std::set} не сортируют, потому что он уже упорядочен. Если нужен другой порядок, обычно делают одно из двух:
\begin{enumerate}
    \item создают новый \texttt{set} с другим компаратором;
    \item копируют элементы в \texttt{std::vector} и сортируют уже его.
\end{enumerate}

\begin{example}
\begin{verbatim}
std::set<int> s = {3, 1, 4};
std::set<int, std::greater<int>> t(s.begin(), s.end());
\end{verbatim}
\end{example}

Это даёт тот же набор элементов, но уже с другим порядком обхода.

\subsection{Как «сортируют» элементы непустого \texttt{std::vector}?}
Для \texttt{std::vector} сортировка выполняется алгоритмом \texttt{std::sort}. Это уже \textit{ручная} сортировка линейного контейнера, в отличие от автоматической упорядоченности \texttt{set} и \texttt{map}.

\begin{example}
\begin{verbatim}
std::vector<int> v = {3, 1, 4, 2};
std::sort(v.begin(), v.end());
\end{verbatim}
\end{example}

Алгоритм \texttt{std::sort} переставляет элементы контейнера так, чтобы они шли в порядке, заданном компаратором.

\subsection{Какие требования накладываются на линейный контейнер, чтобы его можно было сортировать \texttt{std::sort}?}
Для \texttt{std::sort} нужны:
\begin{itemize}
    \item \textbf{random access iterators}, потому что алгоритм активно работает с позицией элементов;
    \item возможность перестановки элементов;
    \item корректный компаратор, задающий строгий слабый порядок.
\end{itemize}

Поэтому \texttt{std::vector} можно сортировать, а \texttt{std::list} -- нельзя, потому что у неё не random access iterators.

\begin{remark}
Сортировка \texttt{std::sort} не требует непрерывной памяти как обязательного свойства, но для \texttt{std::vector} она обычно есть, и это делает контейнер особенно удобным.
\end{remark}

\subsection{Можно ли отсортировать поддиапазон \texttt{std::vector}?}
Да. Алгоритм \texttt{std::sort} работает не только со всем контейнером, но и с любым поддиапазоном, заданным итераторами.

\begin{example}
\begin{verbatim}
std::sort(v.begin() + l, v.begin() + r);
\end{verbatim}
\end{example}

Здесь сортируется только часть вектора на полуинтервале \texttt{[l, r)}.

Это удобно, если нужно упорядочить не весь контейнер, а только его фрагмент.

\subsection{Можно ли упорядочить только часть \texttt{std::set} или \texttt{std::map}?}
Нет, в обычном смысле -- нельзя. Ассоциативный контейнер всегда хранит \textbf{весь} набор элементов в соответствии с одним общим порядком, заданным компаратором.

Нельзя выбрать произвольное подмножество элементов и «отсортировать только его» внутри \texttt{set} или \texttt{map}, потому что:
\begin{itemize}
    \item контейнер уже сам поддерживает порядок;
    \item элементы связаны с внутренним деревом;
    \item ключи в \texttt{set} и \texttt{map} фактически недоступны для произвольного переупорядочивания.
\end{itemize}

Если нужен другой порядок только для части данных, обычно:
\begin{enumerate}
    \item копируют нужные элементы в \texttt{vector};
    \item сортируют \texttt{vector} нужным образом;
    \item либо используют отдельный контейнер с другим компаратором.
\end{enumerate}

\subsection{Можно ли использовать один и тот же компаратор для \texttt{std::set}, \texttt{std::map} и \texttt{std::vector}?}
Да, но только если он подходит по типам сравниваемых объектов.

\begin{itemize}
    \item Для \texttt{std::set<Key, Compare>} и \texttt{std::map<Key, Value, Compare>} компаратор сравнивает ключи типа \texttt{Key}.
    \item Для \texttt{std::vector<T>} компаратор сравнивает элементы типа \texttt{T}.
\end{itemize}

То есть один и тот же \emph{тип} компаратора можно использовать, если он умеет сравнивать нужные значения. Например, один функтор может сравнивать и ключи в \texttt{set}/\texttt{map}, и элементы \texttt{vector}, если все они имеют один и тот же тип и один и тот же смысл порядка.
